\section{结论}
本文从单日确定性计划与波动电价下的连续调度两个层次研究微网购电。以母线侧电量建立供需平衡，通过储能递推连接时段，并以二进制变量约束充放电互斥，形成可求解和复核的混合整数线性规划模型。

问题一得到全天购电量59\,482.699 kWh、购电费用35\,126.949元，给定光伏预测全部消纳，初末储电量均为6\,000 kWh。全天转换损耗3\,940.727 kWh与总能量收支一致，调度在低价和光伏富余时段补充储能，在后续较高电价时段放电，并满足日末恢复要求。

问题四利用已实现联合误差构造三类采购情景，在共享合同和全窗口储能动作的条件下滚动求解，并按实际供需执行、结算。在2025年2月至12月的334日连续回放中，每日合同冻结与允许调整条件的费用分别为1591.9083、1540.2971万元，比同条件点预测基准减少126.2237、46.3193万元，降幅为7.3466\%、2.9194\%。两类紧急购电量分别减少45.7195、28.5367万kWh，实际末态均为6\,000 kWh。

情景机制保留原点预测，通过提高正常合同采购减少高价缺口补电，但合同余电、弃光和储能吞吐增加。因此，总电费降低与预测精度、光伏利用或电池寿命收益不能作简单等同。不同合同条件间的差异还同时包含光伏信息与调整权限的影响。

上述结论以本文效率、费用、信息条件和已查看开发年份为前提。后续需补充问题二、三的独立结果，并结合电池寿命成本、参数扰动及独立年份数据进一步评价策略。有限情景与年末预留均不能保证任意实际供需序列下的可行性。
